%% The comparison table of the introduction.  One row per formulation cited in
%% the text; a dash means the ingredient is not part of that formulation.
\begin{table}[!htbp]
\centering
\small
\caption{Ingredients carried by representative strain-gradient and
flexoelectric formulations. A dash means the ingredient is not part of the
cited formulation; ``splines'' marks $C^1$ isogeometric interpolation. Only
the three last columns distinguish the present work from the one-scale theory
it extends.}
\label{tab:comparison}
\begin{tabular}{@{}lcccccc@{}}
\toprule
formulation & nonlocal or & piezo- & flexo- & frac- & $C^1$ & aniso- \\
 & gradient & electric & electric & tional & & tropic \\
 & & & & memory & & length \\
\midrule
Mindlin, second gradient~\cite{Mindlin1965}          & gradient & -- & -- & -- & -- & -- \\
Eringen, nonlocal elasticity~\cite{Eringen1972Polar} & nonlocal & -- & -- & -- & -- & -- \\
Aifantis~\cite{Aifantis1992}                         & gradient & -- & -- & -- & -- & -- \\
Polizzotto~\cite{Polizzotto2018}                     & gradient & -- & -- & -- & -- & \checkmark \\
Amanatidou and Aravas~\cite{AmanatidouAravas2002}    & gradient & -- & -- & -- & mixed & -- \\
Papanicolopulos and Zervos~\cite{Papanicolopulos2009} & gradient & -- & -- & -- & \checkmark & -- \\
Chen, Yao and Yvonnet~\cite{ChenYaoYvonnet2023}      & gradient & -- & \checkmark & -- & splines & -- \\
Joshan and Santapuri~\cite{JoshanSantapuri2024}      & gradient & -- & \checkmark & -- & -- & -- \\
one-scale theory~\cite{fem1}                         & gradient & -- & -- & -- & \checkmark & \checkmark \\
\textbf{present work} & gradient & \checkmark & \checkmark & \checkmark & \checkmark & \checkmark \\
\bottomrule
\end{tabular}
\end{table}
